\section{The Fall and the Punishment of the Demons}\label{sec:fall}

The angels were created good and in grace, and by one meritorious act
the greater part of them reached the vision of God
(\S\ref{sec:q62}). Some did not. Aquinas considers how the angels
became evil (\latin{quomodo Angeli facti sunt mali}) in two questions:
first as to the evil of fault (\latin{malum culpae}, q.~63), then as to
the evil of punishment (\latin{malum poenae}, q.~64) (\ST{I q.~63 pr.}).
The Church's confession stands behind both.
Lateran~IV teaches that \latin{Diabolus enim et alii daemones a Deo
quidem natura creati sunt boni, sed ipsi per se facti sunt mali. Homo
vero diaboli suggestione peccavit} (\emph{Firmiter}, DS~800), and that
the reprobate will receive \latin{cum diabolo poenam perpetuam}
(DS~801). The Catechism takes up the same words and adds the
consequence: the fall ``consists in the free choice of these created
spirits, who radically and irrevocably rejected God and his reign''
(CCC~391--392).

Scripture gives the fall as event and as figure. ``In his angels he
found wickedness'' (Job 4:18); the devil ``stood not in the truth''
(John 8:44) and ``sinneth from the beginning'' (1 John 3:8); ``God
spared not the angels that sinned, but delivered them, drawn down by
infernal ropes to the lower hell, unto torments, to be reserved unto
judgment'' (2 Pet 2:4; Jude~6); the dragon's tail ``drew the third part
of the stars of heaven'' (Apoc 12:4), and ``his angels were thrown down
with him'' (Apoc 12:9); ``I saw Satan like lightning falling from
heaven'' (Luke 10:18). The Fathers and Aquinas read the oracles against
the king of Babylon (Isa 14:12--15) and the prince of Tyre (Ezek
28:12--17) as spoken to the devil under those persons, and the Vulgate
Job 40--41, where Behemoth is ``the beginning of the ways of God''
(40:14) and Leviathan ``king over all the children of pride'' (41:25),
as spoken of him in figure (Douay throughout). The schools after Aquinas dispute the object and
the moment of the first sin (\S\ref{sec:dq-firstsin};
Appendix~\ref{app:first-sin}). What the demons do to men, and how they
tempt, belongs to q.~114 (\S\ref{sec:assaults}).

\subsection{The Evil of the Angels with Regard to Fault
(\ST{I q.~63})}\label{sec:q63}

Nine articles: whether evil of fault can be in the angels; what kinds
of sin; what the angel desired in sinning; whether some demons are
naturally evil; whether the devil was evil in the first instant of his
creation; whether there was an interval between creation and fall;
whether the highest of those who fell was the highest of all; whether
his sin caused the sin of the others; and whether as many fell as stood
(\ST{I q.~63 pr.}).

\paragraph{The angel could sin (a.~1).} The \emph{sed contra} is Job
4:18. ``An angel or any other rational creature considered in his own
nature, can sin; and to whatever creature it belongs not to sin, such
creature has it as a gift of grace, and not from the condition of
nature.'' Sin is a departure from the rectitude an act ought to have,
and only an act whose rule is the agent's own power cannot depart from
it: were the craftsman's hand itself the rule of the carving, he could
not carve crookedly. Only the divine will is its own rule; every created
will is right only as it is regulated by the divine will, as ``the
soldier's will, according to the will of his commanding officer.'' The
fourth reply fixes how a pure spirit could sin at all. He could not
choose what is evil in itself, for that presupposes error in the
particular, which comes of passion or habit, and the angel had neither.
He could choose what is good in itself ``but not according to proper
measure or rule,'' as a man who prays without heeding the order the
Church has set; such a sin ``does not presuppose ignorance, but merely
absence of consideration of the things which ought to be considered.''
So ``the angel sinned, by seeking his own good, from his own free-will,
insubordinately to the rule of the Divine will'' (ad~4). The natural
love of God as author of nature remains; what could be turned away was
the infused love by which the angel turned to God as the object of
supernatural beatitude (ad~3).

\paragraph{Pride and envy only (a.~2).} The \emph{sed contra} is
Augustine: the devil ``is not a fornicator nor a drunkard, nor anything
of the like sort; yet he is proud and envious'' (\emph{De civitate Dei}
XIV.3); in Augustine's own chapter the devil, having no flesh, ``is
exceedingly proud and envious,'' and on account of it ``is reserved in
chains of darkness to everlasting punishment.'' As to guilt, all sins are in the demons,
since they lead men into all sins; as to affection, only sins that can
belong to a spiritual nature. In spiritual goods there can be no sin in
being drawn to them except by not keeping the rule of the superior,
\latin{et hoc est peccatum superbiae, non subdi superiori in eo quo
debet}. ``Consequently the first sin of the angel can be none other than
pride.'' Envy follows, for to desire a singular excellence is to resist
whatever would make it cease to be singular; ``after the sin of pride,
there followed the evil of envy in the sinning angel, whereby he grieved
over man's good, and also over the Divine excellence.'' Envy here is
taken \latin{non \dots\ pro passione, sed pro voluntate renitente bono
alterius}, for the will resisting another's good (ad~2). The demons delight in the sins of the
flesh not as disposed to them but ``wholly through envy,'' as these
hinder man's good (ad~1). Augustine had put the order in the same terms:
\latin{invidia sequitur superbiam, non praecedit} (\emph{De Genesi ad
litteram} XI.14.18), and Scripture joins the two:
``pride is the beginning of all sin'' (Ecclus 10:15), and ``by the envy
of the devil, death came into the world'' (Wis 2:24).

\paragraph{To be as God (a.~3).} The \emph{sed contra} is Isa 14:13--14,
``I will ascend into heaven \dots\ I will be like the most High.''
``Without doubt the angel sinned by seeking to be as God.'' Not by
equality (\latin{per aequiparantiam}): he knew by natural knowledge that
this was impossible, no passion or habit could lead him to choose the
impossible, and a creature cannot desire a higher grade of nature
without desiring its own destruction, ``just as an ass does not desire
to be a horse.'' He desired it by likeness (\latin{per similitudinem}),
and not the likeness proper to every creature, which is sinless if
sought in due order ``that he may obtain it of God.'' He desired a
likeness proper to God alone: ``by desiring, as his last end of
beatitude, something which he could attain by the virtue of his own
nature, turning his appetite away from supernatural beatitude, which is
attained by God's grace. Or, if he desired as his last end that likeness
of God which is bestowed by grace, he sought to have it by the power of
his own nature'' (\latin{per virtutem suae naturae}). The two readings
``in a manner coincide,'' since on either he sought to have final
beatitude of his own power, ``whereas this is proper to God alone.''
Aquinas finds the same in Anselm, whom he gives as saying that the devil
\latin{appetiit illud ad quod pervenisset si stetisset} (\emph{De casu
diaboli} 4), and adds a consequence: since what is of itself causes what
is of another, ``he sought to have dominion over others.'' Gregory names
him \latin{ille antiquus hostis, qui Deo esse per superbiam similis
concupivit}, the ancient enemy who by pride desired to be like God
(\emph{Hom.~in Ev.} 34.9); the Catechism hears the same rebellion in
``You will be like God'' (CCC~392; Gen 3:5).

\paragraph{Not evil by nature (a.~4).} The \emph{sed contra} is
Dionysius: \latin{Daemones non sunt natura mali} (\emph{De divinis
nominibus} 4). Everything that exists tends naturally to some good,
since it comes from a good principle. A nature inclined to a particular
good may tend accidentally to an evil joined with that good, as fire
consumes; but an intellectual nature is inclined to good in general, to
which no evil is joined, ``hence, since the demons are intellectual
substances, they can in no wise have a natural inclination towards any
evil whatsoever.'' Porphyry held the demons naturally deceitful because
he made them animals with a sensitive nature, and Augustine rebuked him
(ad~1); the slyness of the fox and the fierceness of the dog are no evil
in them, as Dionysius observes (ad~3). The Church has taught the article
against the Manichees and Priscillianists: the First Council of Braga
anathematizes whoever says \latin{diabolum non fuisse prius (bonum)
angelum a Deo factum, nec Dei opificium fuisse naturam eius}, but that
he arose from chaos and darkness and is himself \latin{principium atque
substantiam mali} (can.~7, DS~457), and Lateran~IV confesses the
demons \latin{natura creati \dots\ boni} (DS~800). Damascene had said
the same in the East: the devil ``was not made wicked in nature but was
good,'' and ``of his free choice was changed'' (\emph{De fide orthodoxa}
II.4).

\paragraph{Not evil in the first instant (a.~5).} The \emph{sed contra}
is Gen 1:31, ``God saw all the things that he had made, and they were
very good.'' Some held the devil wicked from the first instant, not by
nature but by his own will refusing righteousness as soon as he was
made; Augustine grants that such a view is not Manichaean (\emph{De
civitate Dei} XI.13). But Isa 14:12, ``O Lucifer, who
didst rise in the morning,'' and Ezek 28:13, ``Thou wast in the
pleasures of the paradise of God,'' stand against it, and so
\latin{a magistris haec opinio tanquam erronea rationabiliter reprobata
est}. Others said the angel could have sinned in the first instant but
did not; the argument some brought against them, that two successive
changes cannot terminate in one instant, fails for instantaneous
changes, since the air is lit by the moon in the same instant the moon
is lit by the sun. Aquinas's own reason is from the cause: an operation
that begins with a thing's existence comes from the agent that gave it
its nature, as a leg defective from birth ``begins at once to limp''; but
God ``cannot be the cause of sin.'' ``He was a murderer from the
beginning'' (John 8:44) means, with Augustine, ``from the beginning of
sin'' (ad~1; \emph{De civitate Dei} XI.15). All the angels, created in
grace, merited in the first instant; ``some of them at once placed an
impediment to their beatitude, thereby destroying their preceding
merit'' (ad~4). \emph{De malo} q.~16 a.~4 reaches the same conclusion by
another way, noted below.

\paragraph{The fall immediately after the first instant (a.~6).} The
\emph{sed contra} is John 8:44 with Augustine: ``he was in the truth,
but did not remain in it.'' There is a twofold opinion, but \latin{
probabilior, et sanctorum dictis magis consona est, quod statim post
primum instans suae creationis Diabolus peccaverit}. This follows if
the angel elicited an act of free will in the first instant and was
created in grace (\S\ref{sec:q62}): having merited by one act, he would
have been beatified at once after that instant, had he not sinned. ``If,
however, it be contended that the angel was not created in grace, or
that he could not elicit an act of free-will in the first instant, then
there is nothing to prevent some interval being interposed between his
creation and fall.'' The ``walking'' of Ezek 28:15 figures the
instantaneous movement of free will (ad~1). Man had a delay before his
sin; the angel's free will is inflexible after choice, so that had he
not sinned at once he would have been confirmed in good (ad~3). Angelic
time is the succession of the angels' acts. In the first instant each
knew himself by evening knowledge, and this act was good in all; in the
second, some turned to the Word in morning knowledge, while others,
absorbed in themselves, ``became night, `swelling up with pride'\,''
(\latin{facti sunt nox, per superbiam intumescentes}), as Aquinas cites
Augustine (ad~4; on evening and morning knowledge, \S\ref{sec:q58}).

\paragraph{The highest who sinned was probably the highest of all
(a.~7).} The \emph{sed contra} is Gregory: the first angel, ``set over
all the hosts of angels, surpassed them in brightness, and was by
comparison the most illustrious among them'' (\latin{dum cunctis
agminibus angelorum praelatus est, ex eorum comparatione clarior fuit},
\emph{Hom.~in Ev.} 34.7). Two things weigh on the
question, proneness to sin and the motive of sinning. By proneness the
higher angels were less likely to fall, and Damascene accordingly holds
that the chief of the fallen was set over the earthly order (\emph{De
fide orth.} II.4); nor is this opinion ``to be rejected as contrary to
faith,'' since God governs bodies through the angels. By motive the
higher were more exposed, for pride's motive is excellence. Since the
angels' sin came not of any proneness but ``of free choice alone,'' the
argument from motive weighs more, \latin{et hoc videtur probabilius};
``yet this must not be prejudicial to the other view.'' Ezekiel calls
him ``a cherub'' (Ezek 28:14) and not a seraph, because Cherubim is
named from fullness of knowledge, which is compatible with mortal sin,
and Seraphim from the heat of charity, which is not (ad~1). Gregory reads
Job 40:14 the same way: God ``created him first, whom He made more
eminent than the other Angels,'' and ``when sinning, he was condemned
without pardon, because he had been created great beyond comparison''
(\emph{Moralia} XXXII.23.47).

\paragraph{The first sinner drew the others (a.~8).} The \emph{sed
contra} is Apoc 12:4. \latin{Peccatum primi Angeli fuit aliis causa
peccandi, non quidem cogens, sed quadam quasi exhortatione inducens.}
The sign is that all the demons are subject to him: ``Depart from me,
you cursed, into everlasting fire, which was prepared for the devil and
his angels'' (Matt 25:41), for divine justice subjects whoever consents
to another's suggestion to him in punishment: ``by whom a man is
overcome, of the same also he is the slave'' (2 Pet 2:19). Exhortation
and consent need no time among angels, so all could sin in one instant
and one sin still be the cause of the rest (ad~1). The proud would
rather serve a superior, but chose to serve one of their own nature so
as to reach beatitude by natural powers (ad~2). The highest, moved with
his whole greater energy, ``became the greater in malice'' (ad~3).
Damascene gives the same picture: ``an innumerable host of angels
subject to him were torn away and followed him'' (\emph{De fide orth.}
II.4).

\paragraph{More stood than fell (a.~9).} The \emph{sed contra} is
4 Kgs 6:16, ``there are more with us than with them,'' expounded of the
good angels and the wicked spirits. \latin{Plures Angeli permanserunt
quam peccaverunt}, since sin is against natural inclination, and what
is against nature happens in the fewer cases. Among men evil is in the
many because sense goods are known to most and the good of reason to
few; the angels have only an intellectual nature (ad~1). The third of
the stars (Apoc 12:4) is thus the figure of a minority. If the chief
devil was of the lowest order, only some of that order fell; if of the
highest, ``it is probable that some fell of every order,'' as men are
taken into every order to repair the ruin. Scripture never gives the
demons the names Seraphim and Thrones, which are drawn from charity and
God's indwelling, but does give them Cherubim, Powers, and
Principalities (ad~3; on the orders, Appendix~\ref{app:orders}).
Augustine holds the ruin to be filled from men, and confesses: ``We do
not know the number either of the saints or of the devils''
(\emph{Enchiridion} 29).

\subsubsection*{The Parallel in \emph{De malo} q.~16}

The sixteenth question of \emph{De malo}, on the demons, is the
disputed form of the same matter and extends it to the demons' bodies,
knowledge, and power over bodies and minds. Its first five articles run
with qq.~63--64 of the \emph{Summa}; the rest answer to the treatise's
earlier questions and to q.~114.

\begin{center}\small
\begin{tabularx}{\linewidth}{@{}l>{\raggedright\arraybackslash}p{0.27\linewidth}>{\raggedright\arraybackslash}X@{}}
\toprule
\emph{De malo} & Question & Relation to the \emph{Summa} \\
\midrule
q.~16 a.~1 & Bodies naturally united & Agrees with q.~51 a.~1; the point \latin{non multum refert ad fidei Christianae doctrinam} \\
a.~2 & Evil by nature or will & Agrees with q.~63 a.~4 \\
a.~3 & Sin as desire of equality & Agrees with q.~63 a.~3; sin in the supernatural order \\
a.~4 & First instant & Same conclusion as q.~63 a.~5; different reason \\
a.~5 & Free will after the fall & Agrees with q.~64 a.~2; \latin{vis primae electionis} \\
a.~6 & Deception and false opinion & Refines q.~64 a.~1 \\
aa.~7--8 & Futures; thoughts of hearts & Agrees with q.~57 aa.~3--4 \\
aa.~9--12 & Bodies, motion, sense, intellect & Answers to qq.~110--111, 114 \\
\bottomrule
\end{tabularx}
\end{center}

Three articles differ in emphasis or argument. On the sin itself
(a.~3), the demon willed final beatitude \latin{per virtutem suae
naturae; non tamen sine Deo in naturam operante, sed sine Deo gratiam
conferente}: the sin lies not in a natural good but in refusing to have
the supernatural end as gift. On the first instant (a.~4), the
disputation names the opinion of \latin{quidam moderni} and records
that it \latin{reprobata fuit ab omnibus magistris tunc Parisiis
legentibus}; but it sets aside the argument from the nature as God
created it, since the malice of fault is founded in a good nature as in
its subject (\latin{malitia culpae non repugnat bonitati naturae, sed in
ea fundatur sicut in subiecto}), and argues instead from the order of
the angel's acts. In the
first instant the angel turns to natural knowledge of himself, in which
he cannot sin; only afterward to what is above nature, or away from it;
so that \latin{in primo instanti suae creationis non fuit neque beatus
per conversionem perfectam in Deum, neque peccator per aversionem ab
ipso}. On obstinacy (a.~5), the \emph{De malo} joins to the intrinsic
immobility of the angelic will the extrinsic end of the angel's state of
pilgrimage, and concludes that the demons sin in all they choose,
\latin{quia in omni eorum electione permanet vis primae electionis
ipsorum}. The disputation's argument on the first instant returns in the
schools after Aquinas (\S\ref{sec:dq-firstsin}).

\angeldossier{\ST{I q.~63 aa.~1--9}; parallels \emph{In II Sent.} dd.~5--6;
\emph{De malo} q.~16 aa.~2--4.}
 {Defined: the devil and the other demons were created good in nature by
 God and became evil of themselves (Lateran~IV, \emph{Firmiter},
 DS~800); the devil was first a good angel made by God and is no
 principle or substance of evil (Braga~I, can.~7, DS~457) (aa.~4--5).
 Church teaching: the fall was a free, radical, and irrevocable
 rejection of God, reflected in ``You will be like God'' (CCC~391--392,
 414). Common teaching: the angel could sin by nature and is impeccable
 only by grace (a.~1); the first sin was pride, and envy its consequence
 (a.~2); the angel sought to be as God by likeness, not by equality
 (a.~3); the chief of the fallen induced the rest and rules them (a.~8);
 more angels stood than fell (a.~9). Thomist position, as the more
 probable: the object of the pride was beatitude by the angel's own
 power (a.~3); the fall came immediately after the first instant (a.~6);
 the first sinner was the highest of all the angels (a.~7). Disputed in
 the schools: the object and the moment of the first sin
 (\S\ref{sec:dq-firstsin}; Appendix~\ref{app:first-sin}).}
 {Job 4:18 (\emph{sed contra} of a.~1); Isa 14:12--14 and Ezek 28:13--15
 (aa.~3, 5, 6); Gen 1:31 (a.~5); John 8:44 (aa.~5--6); Apoc 12:4
 (aa.~8--9); Matt 25:41 and 2 Pet 2:19 (a.~8); 4 Kgs 6:16 (a.~9), all
 Douay. Augustine, \emph{De civitate Dei} XI.13, XI.15, XIV.3, and
 \emph{Enchiridion} 28--29; \emph{De Genesi ad litteram} XI.14.18 and
 XI.23.30 (\latin{superbia tumidus}); \emph{De Gen.\ ad litt.} IV.24
 (a.~6 ad~4) and \emph{Quaest.\ Vet.\ et Novi Test.} 113 (a.~3). Anselm,
 \emph{De casu diaboli} 4 (a.~3). Gregory, \emph{Hom.~in Ev.} 34.7, 34.9,
 and \emph{Moralia} XXXII.23.47. Damascene, \emph{De fide orth.} II.4.
 Dionysius, \emph{De div.\ nom.} 4 (a.~4). Lateran~IV (DS~800); Braga~I
 (DS~457); CCC~391--392, 414.}
 {The opinion that the devil was evil by his own will from the first
 instant, which Augustine does not count Manichaean and the Paris
 masters reproved (a.~5; \emph{De malo} q.~16 a.~4). Damascene's view
 that the chief of the fallen presided over the earthly order and that
 the fallen were of the lower grade (a.~7), which Aquinas declines to
 call contrary to faith. The opinion that the angels were created
 without grace, on which an interval between creation and fall is
 possible (a.~6; \S\ref{sec:dq-grace}).}

\subsection{The Punishment of the Demons (\ST{I q.~64})}\label{sec:q64}

Four articles: whether the demons' intellect is darkened by privation of
the knowledge of all truth; whether their will is obstinate in evil;
whether there is sorrow in them; and whether this atmosphere is their
place of punishment (\ST{I q.~64 pr.}).

\paragraph{Knowledge darkened, nature untouched (a.~1).} The \emph{sed
contra} is Dionysius: the gifts given to the demons \latin{nequaquam
mutata esse dicimus, sed sunt integra et splendidissima} (\emph{De div.\
nom.} 4); in Parker's English they ``exist, and are complete, and all
luminous, although the demons themselves do not see, through having
blunted their powers of seeing good'' (\emph{DN} 4.23). Knowledge of
truth is natural or of grace, and that of grace is either speculative,
the revelation of divine secrets, or affective, \latin{affectiva,
producens amorem Dei}, producing love for God and belonging to the gift
of wisdom. The first ``was neither taken away nor lessened,'' because
nothing can be withdrawn from a simple nature ``as a man is punished by
being deprived of a hand or a foot.'' The second is lessened, for only
so much is revealed to them as is needed, by angels or ``through some
temporal workings of Divine power,'' as Aquinas cites Augustine (\emph{De
civitate Dei} IX.21). ``But of the third knowledge, as likewise of
charity, they are utterly deprived.'' Their knowledge of creatures,
not referred to God, is not evening but night: \latin{non dicitur
vespertina, sed nocturna} (ad~3). They knew the mystery of the
Incarnation less fully; ``for had they fully and certainly known that He
was the Son of God and the effect of His passion, they would never have
procured the crucifixion of the Lord of glory'' (ad~4). \emph{De malo}
q.~16 a.~6 adds a precision: no false opinion can be in them about what
falls under their natural knowledge, but there can be about what exceeds
it.

\paragraph{The will obstinate in evil (a.~2).} The \emph{sed contra} is
Ps 73:23, ``the pride of them that hate thee ascendeth continually,''
understood of the demons. Aquinas reports Origen's opinion that every
created will but Christ's soul can be inclined to good and to evil
(\emph{Peri Archon} I.6). That view takes from angels and saints the
stability of beatitude, and contradicts Scripture's ``everlasting
punishment'' and ``life everlasting'' (Matt 25:46); \latin{unde haec
positio est tanquam erronea reputanda; et tenendum est firmiter, secundum
fidem Catholicam, quod et voluntas bonorum Angelorum confirmata est in
bono, et voluntas Daemonum obstinata est in malo}. The cause is not the
gravity of the sin but ``the condition of their nature or state,'' for,
as Aquinas cites Damascene, ``death is to men, what the fall is to the
angels.'' Damascene's own sentence runs: ``what in the case of man is
death is a fall in the case of angels. For after the fall there is no
possibility of repentance for them, just as after death there is for men
no repentance'' (\emph{De fide orth.} II.4). The appetite follows the
apprehension. Man
reasons from one thing to another and holds his objects movably; the
angel \latin{apprehendit immobiliter per intellectum}, as we hold first
principles, and his will \latin{adhaeret fixe et immobiliter}. Hence the
angel's free will \latin{est flexibile ad utrumque oppositum ante
electionem, sed non post}. God's mercy frees the penitent, but those
``not capable of repenting, cling immovably to sin'' (ad~2). The first
sin remains as desire, though not as the belief that its end can be
had (ad~3). When a demon confesses Christ or believes and trembles, he
does good but not well, ``compelled by the evidence of things''
(ad~5; Mark 1:24; James 2:19). Gregory: ``the heart of the ancient enemy
will be hardened as a stone, because it will never be softened by any
penitence of conversion'' (\emph{Moralia} XXXIV.6.11). The Church has
condemned the contrary. The anathematisms against Origen published in
the synod of Constantinople of 543 strike whoever holds that Christ will
be crucified
\latin{pro daemonibus, sicuti et pro hominibus} (can.~7, DS~409), or
that the punishment of the demons is temporary and will end, \latin{sive
restitutionem et redintegrationem esse (fore) daemonum aut impiorum
hominum} (can.~9, DS~411). The Catechism speaks from Damascene: ``It is
the irrevocable character of their choice, and not a defect in the
infinite divine mercy, that makes the angels' sin unforgivable''
(CCC~393).

\paragraph{Sorrow in the demons (a.~3).} The \emph{sed contra} is Apoc
18:7, ``as much as she hath glorified herself \dots\ so much torment and
sorrow give ye to her'': much more is the devil punished with sorrow,
``because he especially glorified himself.'' Fear, sorrow, and joy as
passions belong to the sensitive appetite and cannot be in the demons;
as simple acts of the will they can. Sorrow so taken ``is nothing else
than the resistance of the will to what is, or to what is not''
(\latin{renisus voluntatis}), and the demons would have many things not
be which are, and others be which are not, wishing out of envy the
damnation of the saved; they are deprived of the happiness they desire
naturally, and their wicked will is curbed in many ways. They can
rejoice in one thing and grieve in another (ad~1), and ``the devils also
believe and tremble'' (ad~2; James 2:19). They are sorry for the
punishment, which witnesses to the goodness of nature, but not for the
sin, which would witness to a good will; ``since the demon has a
perverse and obstinate will, he is not sorry for the evil of sin''
(ad~3).

\paragraph{Hell and the dark air (a.~4).} The \emph{sed contra} is
Augustine: ``the darksome atmosphere is as a prison to the demons until
the judgment day'' (\latin{ista caliginosa tenere permissa sunt, qui eis
pro suo genere quidam quasi carcer est, usque ad tempus iudicii},
\emph{De Genesi ad litteram} III.10.15). The angels
stand by nature between God and men, and providence procures the good
of the lower through the higher: directly, through the good angels, and
indirectly, as one who is assailed is exercised by the struggle, which
fittingly is done through the wicked spirits, ``lest they should cease
to be of service in the natural order.'' So \latin{Daemonibus duplex
locus poenalis debetur}: hell, by reason of their fault, and \latin{
caliginosus aer}, by reason of the exercise of men. This lasts as long
as the work of salvation, until the judgment; ``some of them are even
now in hell, to torment those whom they have led astray''; after the
judgment all the wicked, men and angels, will be in hell. A place
punishes a spirit by grieving the will, not by changing the nature
(ad~1). While in the air they are not bound in the fire, but their
punishment is not less, ``because they know that such confinement is
their due,'' and a gloss on James 3:6 says ``they carry fire of hell
with them wherever they go'' (ad~3). When they ``besought him that he
would not command them to go into the abyss'' (Luke 8:31), they asked
not to be shut out from the place where they could harm men (ad~3; Mark
5:10). Scripture names both places: ``the lower hell'' (2 Pet 2:4) and
``the spirits of wickedness in the high places'' (Eph 6:12). Augustine,
leaving other questions of the fall open, holds this without doubt:
\latin{peccatores angelos minime dubitemus detrusos tamquam in carcerem
caliginis huius aeriae circa terras, secundum apostolicam fidem, in
iudicio puniendos servari}, the sinning angels thrust into the prison of
this dark air about the earth and kept for punishment at the judgment
(\emph{De Gen.\ ad litt.} XI.26.33). The demons' present liberty to
tempt is treated with q.~114 (\S\ref{sec:assaults}); the Catechism
states its limit: ``The power of Satan is, nonetheless, not
infinite. He is only a creature'' (CCC~395).

\angeldossier{\ST{I q.~64 aa.~1--4}; parallels \emph{In II Sent.}
dd.~6--7; \emph{De malo} q.~16 aa.~1, 5--6.}
 {Defined: the reprobate will suffer perpetual punishment with the devil
 (Lateran~IV, DS~801); the anathematisms against Origen had already
 struck the demons' restoration (Constantinople 543, can.~9, DS~411) and
 a future crucifixion of Christ for them (can.~7, DS~409). Church
 teaching: the angels' choice is irrevocable and their sin unforgivable,
 not by defect of mercy (CCC~393, 414); Satan's power is a creature's and
 not infinite (CCC~395). Common teaching: the will of the demons is
 obstinate in evil, which Aquinas holds as of Catholic faith (a.~2);
 their natural gifts remain entire, their graced speculative knowledge is
 lessened, and their affective knowledge and charity are lost (a.~1);
 sorrow is in them as an act of will (a.~3); the dark air is a second
 place of their punishment until the judgment (a.~4, with Augustine,
 \emph{De Gen.\ ad litt.} III.10.15, XI.26.33). Thomist position:
 obstinacy is caused by the immobility of the angelic apprehension and
 will, as death fixes man's (a.~2, with Damascene, \emph{De fide orth.}
 II.4).}
 {Ps 73:23 and Matt 25:46 (a.~2); Apoc 18:7, James 2:19, Job 41:24
 (a.~3); Luke 8:31, Mark 5:10, 2 Pet 2:4, Eph 6:12 (a.~4), all Douay.
 Dionysius, \emph{De div.\ nom.} 4.23 (\emph{sed contra} of a.~1).
 Augustine, \emph{De Gen.\ ad litt.} III.10.15 (\emph{sed contra} of
 a.~4) and XI.26.33; \emph{De civitate Dei} IX.21, XIX.13, and
 \emph{De Gen.\ contra Manich.} II.17 (aa.~1, 3). Damascene, \emph{De
 fide orth.} II.4 (a.~2). Gregory, \emph{Moralia} XXXIV.6.11. Origen,
 \emph{De principiis} I.6 (a.~2); gloss on James 3:6 (a.~4 ad~3).
 Lateran~IV (DS~801); anathematisms against Origen (DS~409, 411);
 CCC~393--395, 414.}
 {Origen's opinion that every created will save Christ's soul remains
 flexible to good and evil, and that the devil's orders may be
 converted, which Origen himself leaves to the reader (\emph{De princ.}
 I.6.3) and Aquinas judges erroneous (a.~2). The opinion that the pain
 of sense of demons and souls is postponed until the judgment day, which
 Aquinas calls erroneous as to souls; as to demons, ``it is better to
 say that the same judgment is passed upon wicked souls and wicked
 angels'' (a.~4 ad~3).}
